% Pursuing Stacks (cotes 134-2 to 134-8): archive cote 134-5, site page 110 (PDF page 111), typescript page 297.
% typescript. Draft compared: ps5.tex:284-345.
% AI-transcribed by Claude Opus 5.5 (claude-opus-5-5, adaptive thinking, effort high), Message Batches API, batch ps2 (msgbatch_01NVwr1Yuo8qJf5rykeGBXPK), 2026-11-12; not yet checked by eye.
\scanpage{134-5/110}
\mspage{297}
equal to a chosen final object e of A,  in order to get a functor
\[ I \longmapsto I_A \ : \ \text{(Sets)} \longrightarrow A \ , \]
where $I_A$  is the ``constant'' object of A with value I,  namely a sum of I
copies of e (sometimes also written  Ixe ).  Using this functor,  we get
a functor
\[ \hand{\Delta}^{\wedge} = \underline{\mathrm{Hom}}(\hand{\Delta}^{\circ}, \text{(Sets)}) \longrightarrow M = \underline{\mathrm{Hom}}(\hand{\Delta}^{\unsure{\circ}} , A) \ , \]
which I denote by
\[ K \longmapsto K_A \qquad , \]
associating to any ss set the corresponding ``constant'' \inserted[typed]{(relative to A)} ss object of A.
On the other hand,  because A admits finite products,  so does M,  which
enables us to make use of the homotopy notions developped in sections 51
etc.  Thus if
\[ \underline{I} = (I, \hand{\delta}_0, \hand{\delta}_1) \]
is any interval in $\hand{\Delta}^{\wedge}$,  considering the corresponding ``\inserted{A-}constant''  inter-%
val  $\underline{I}_A$  in M,  we get homotopy notions in M,  \unsure{which} we may refer to as
$\underline{I}$ - \uline{homotopy} (dropping the subscript A).  They can all be deduced from the
elementary \inserted{$\underline{I}$-} homotopy \xout{relation} between maps \inserted{in M},  which  is expressed in the
known way,  in terms of a map in  M
\[ h\colon \underline{I}_A \ x \ X \longrightarrow Y \quad , \]
where X and Y are the source and target in M of the two considered maps,
between which we want to find an \xout{homotopy} elementary $\underline{I}$-homotopy h.  This
map h  decomposes componentwise into
\[ (*\ ) \qquad h_n : \hand{(}I_n\hand{)}_A \ x \ X_n \longrightarrow Y_n \quad , \]
and each $h_n$ can be interpreted , in view of the definition of $I_n$ ,  as a
\inserted[typed]{map}\xout{map}
\[ (*\ ') \qquad h'_n \ : \ I_n \longrightarrow \mathrm{Hom}_A(X_n,Y_n) \quad , \]
provided we assume that taking products in A is distributive with respect
to the sums we are taking,  whence
\[ I_{n\ A} \ x \ X_n \simeq (I_n)_{X_n} \ = \ \text{direct sum of } I_n \text{ copies of } \inserted[typed]{$X_n$.} \]
Now to give  h ,  or equivalently a sequence of maps $h_n$ in (*)  ``functo-%
rial in n for variable n'' ,  amounts to giving a sequence of maps (* '),
satisfying a corresponding compatibility relation  for variable n.  The
point of course is that  ( for any I in $\hand{\Delta}^{\wedge}$ and  X,Y  in M)  \uline{the set of}
\uline{data} (*' )  \uline{plus the compatibility condition make\xout{s} a sense},  \uline{formally},
\uline{independently of}  \uline{any exactness assumptions on A}.  Thus it can be taken

\margin{110}
